\documentclass[stu]{apa7}

\usepackage[spanish]{babel}
\usepackage[utf8]{inputenc}
\usepackage{hyperref}
\usepackage{enumitem}
\usepackage{geometry}
\usepackage{times}
\usepackage{array}
\usepackage{longtable}
\usepackage{booktabs}
\usepackage{colortbl}
\usepackage{xcolor}
\usepackage{tabularx}
\usepackage{multirow}

\geometry{margin=1in}

\title{Actividad 7 --- Fase Final de Implementación \\ Sistema de información web contable sobre recibos de retención y movimientos de caja de la empresa minera MINCH SRL.}
\author{...}
\affiliation{...}
\course{...}
\professor{...}
\duedate{...}

\setlength{\parindent}{0.5in}

\newcolumntype{L}[1]{>{\raggedright\arraybackslash}p{#1}}
\newcolumntype{C}[1]{>{\centering\arraybackslash}p{#1}}

\begin{document}

\maketitle

\section*{Enfoque: Scrum --- Sprint Final y Preparación del Release}

En esta etapa, se realiza un \textbf{Sprint de Cierre} para estabilizar el sistema SIC-MINCH, completando los módulos pendientes identificados en la Actividad 4. El objetivo es refinar la UI/UX, optimizar consultas, implementar seguridad definitiva (roles y permisos) y preparar el despliegue final en producción.

\section{Revisión del Product Backlog y Priorización Final}

Al finalizar el Sprint 1, el equipo completó 9 historias (40 SP) correspondientes a los módulos de administración (usuarios, roles, permisos, proveedores, cuentas bancarias, impuestos) y retenciones (registro, cálculo automático, numeración, recibo PDF). Quedan pendientes 13 historias (68 SP) que deben completarse en el Sprint Final.

\subsection{Historias completadas en Sprint 1}

\begin{longtable}{|C{1cm}|L{4.5cm}|C{2cm}|C{1.5cm}|}
\hline
\rowcolor[gray]{0.9}
\textbf{ID} & \textbf{Historia (resumen)} & \textbf{Épica} & \textbf{SP} \\
\hline
HU-21 & Gestionar cuentas bancarias (CRUD) & EP-06 & 3 \\
\hline
HU-22 & Gestionar impuestos/tasas (CRUD) & EP-06 & 3 \\
\hline
HU-19 & Gestionar proveedores (CRUD) & EP-06 & 3 \\
\hline
HU-17 & Gestionar usuarios (CRUD) & EP-06 & 5 \\
\hline
HU-18 & Gestionar roles y permisos & EP-06 & 5 \\
\hline
HU-01 & Registrar retención con cálculo automático & EP-01 & 8 \\
\hline
HU-02 & Calcular descuentos automáticamente & EP-01 & 5 \\
\hline
HU-03 & Numerar retención automáticamente & EP-01 & 3 \\
\hline
HU-04 & Generar recibo de retención PDF & EP-01 & 5 \\
\hline
\rowcolor[gray]{0.9}
\textbf{Total Sprint 1} & & & \textbf{40} \\
\hline
\end{longtable}

\subsection{Historias restantes (priorizadas por valor de negocio)}

\begin{longtable}{|C{0.8cm}|L{5cm}|C{2cm}|C{1.5cm}|}
\hline
\rowcolor[gray]{0.9}
\textbf{ID} & \textbf{Historia (resumen)} & \textbf{Prioridad (MoSCoW)} & \textbf{SP} \\
\hline
HU-06 & Registrar movimiento bancario (débito/crédito) asociado a una cuenta & Indispensable & 8 \\
\hline
HU-08 & Ver libro de bancos de un mes con saldo corriente en cada fila & Indispensable & 5 \\
\hline
HU-10 & Registrar movimiento de caja (ingreso/egreso) & Indispensable & 8 \\
\hline
HU-13 & Consultar estado de cuenta de un proveedor con movimientos y saldo & Indispensable & 8 \\
\hline
HU-15 & Generar reportes con filtros por fecha y tipo (PDF/Excel) & Indispensable & 8 \\
\hline
HU-07 & Numerar transacción bancaria por año fiscal (octubre--septiembre) & Indispensable & 3 \\
\hline
HU-11 & Numerar movimiento de caja por mes calendario & Indispensable & 3 \\
\hline
HU-05 & Exportar retenciones de un mes a Excel & Importante & 5 \\
\hline
HU-09 & Generar recibo PDF de banco & Importante & 3 \\
\hline
HU-12 & Generar recibo PDF de caja & Importante & 3 \\
\hline
HU-14 & Generar PDF del estado de cuenta & Importante & 3 \\
\hline
HU-16 & Dashboard ejecutivo con KPIs y gráficos & Importante & 8 \\
\hline
HU-20 & Gestionar clientes (CRUD) & Deseable & 3 \\
\hline
\end{longtable}

\subsection{Planificación del Sprint Final}

\textbf{Objetivo del Sprint Final:} \enquote{Completar los módulos de libro de bancos, caja, estados de cuenta y reportes, integrando todas las funcionalidades pendientes y preparando el sistema para producción.}

\textbf{Capacidad del equipo:} 45 SP (3 desarrolladores $\times$ 10 días hábiles $\times$ 0.75 factor enfoque $\approx$ 22.5 días-hombre ajustados).

\textbf{Historias seleccionadas para el Sprint Final:}

\begin{longtable}{|C{0.8cm}|L{4.5cm}|C{2cm}|C{1.5cm}|C{1.5cm}|}
\hline
\rowcolor[gray]{0.9}
\textbf{ID} & \textbf{Historia (resumen)} & \textbf{Épica} & \textbf{SP} & \textbf{Responsable} \\
\hline
HU-06 & Registrar movimiento bancario & EP-02 & 8 & Técnico 1 \\
\hline
HU-07 & Numerar transacción bancaria & EP-02 & 3 & Técnico 1 \\
\hline
HU-08 & Ver libro de bancos con saldo corriente & EP-02 & 5 & Técnico 2 \\
\hline
HU-10 & Registrar movimiento de caja & EP-03 & 8 & Técnico 1 \\
\hline
HU-11 & Numerar movimiento de caja & EP-03 & 3 & Técnico 1 \\
\hline
HU-13 & Consultar estado de cuenta & EP-04 & 8 & Técnico 2 \\
\hline
HU-15 & Generar reportes con filtros & EP-05 & 8 & Técnico 2 \\
\hline
\rowcolor[gray]{0.9}
\textbf{Total} & \textbf{7 historias} & & \textbf{43} & \\
\hline
\end{longtable}

Las siguientes historias quedan fuera del Sprint Final y se recomienda abordarlas como mejoras post-release (sprint posterior o mantenimiento evolutivo):

\begin{itemize}
    \item HU-05 (Exportar retenciones a Excel) --- 5 SP
    \item HU-09 (Recibo PDF de banco) --- 3 SP
    \item HU-12 (Recibo PDF de caja) --- 3 SP
    \item HU-14 (PDF estado de cuenta) --- 3 SP
    \item HU-16 (Dashboard ejecutivo) --- 8 SP
    \item HU-20 (Gestionar clientes) --- 3 SP
\end{itemize}

\section{Plan de Actividades del Sprint Final}

\subsection{Descomposición en tareas}

A continuación se presenta el desglose de tareas para cada historia seleccionada en el Sprint Final.

\subsubsection{HU-06: Registrar movimiento bancario (8 SP)}

\begin{enumerate}
    \item Crear tabla transactions (migración) con campos: id, number, date, amount, payment\_type (T/CH), number\_check, account\_id, movement\_id.
    \item Implementar modelo Transaction con relaciones: belongsTo Movement, belongsTo Account.
    \item Implementar boot event creating con numeración por año fiscal y lockForUpdate().
    \item Desarrollar Livewire component BankMovementComponentForm con formulario de registro.
    \item Integrar búsqueda de proveedores y selección de cuenta bancaria.
    \item Validar: tipo de pago obligatorio, número de cheque obligatorio si payment\_type = CH.
    \item Pruebas unitarias de creación de movimiento bancario.
\end{enumerate}

\subsubsection{HU-07: Numerar transacción bancaria por año fiscal (3 SP)}

\begin{enumerate}
    \item Implementar lógica de año fiscal: octubre a septiembre.
    \item Consultar último número del mismo tipo (D/C) en el año fiscal con lockForUpdate().
    \item Asignar número secuencial: D-000001, C-000001.
    \item Pruebas unitarias de numeración concurrente y cambio de año fiscal.
\end{enumerate}

\subsubsection{HU-08: Ver libro de bancos con saldo corriente (5 SP)}

\begin{enumerate}
    \item Desarrollar Livewire component BankBookComponent con vista mensual por cuenta.
    \item Implementar consulta SQL con función de ventana SUM OVER para saldo corriente.
    \item Columnas: fecha, descripción, débito, crédito, saldo.
    \item Filtros: mes, cuenta bancaria, proveedor.
    \item Pruebas de integración con datos de prueba.
\end{enumerate}

\subsubsection{HU-10: Registrar movimiento de caja (8 SP)}

\begin{enumerate}
    \item Crear tabla boxes (migración) con campos: id, number, movement\_id.
    \item Implementar modelo Box con boot event creating para numeración por mes.
    \item Desarrollar Livewire component CashMovementComponentForm con formulario de registro.
    \item Integrar búsqueda de proveedores.
    \item Calcular saldo anterior si es el primer movimiento del mes.
    \item Validar: tipo (D/C), monto, fecha, proveedor obligatorios.
    \item Pruebas unitarias de creación de movimiento de caja.
\end{enumerate}

\subsubsection{HU-11: Numerar movimiento de caja por mes calendario (3 SP)}

\begin{enumerate}
    \item Implementar numeración por tipo (D/C) dentro del mes calendario.
    \item Consultar último número del mes con lockForUpdate().
    \item Formato: DOC-000001.
    \item Pruebas unitarias de secuencialidad mensual.
\end{enumerate}

\subsubsection{HU-13: Consultar estado de cuenta (8 SP)}

\begin{enumerate}
    \item Desarrollar Livewire component AccountStatementComponent con pestañas (Proveedores/Clientes).
    \item Listar titulares con saldo actual (débito $-$ crédito) usando AccountStatementService.
    \item Vista de detalle con movimientos paginados y saldo corriente.
    \item Mostrar totales: suma débitos, suma créditos, saldo final.
    \item Filtro por rango de fechas y tipo de persona.
    \item Pruebas de integración con el servicio de estados de cuenta.
\end{enumerate}

\subsubsection{HU-15: Generar reportes con filtros (8 SP)}

\begin{enumerate}
    \item Desarrollar Livewire component ReportComponent con selector de tipo (retenciones/caja/banco).
    \item Filtros: rango de fechas, tipo de movimiento, cuenta (si aplica).
    \item Tabla de resultados con saldo corriente.
    \item Implementar exportación a PDF con mPDF.
    \item Implementar exportación a Excel con Maatwebsite.
    \item Pruebas de exportación con diferentes combinaciones de filtros.
\end{enumerate}

\subsection{Planificación de tareas del Sprint Final}

\begin{longtable}{|C{0.8cm}|L{3.5cm}|L{2.5cm}|C{2cm}|C{2cm}|C{2cm}|}
\hline
\rowcolor[gray]{0.9}
\textbf{HU} & \textbf{Tarea} & \textbf{Responsable} & \textbf{Estado (inicio)} & \textbf{Estado (final)} & \textbf{Fecha entrega} \\
\hline
HU-06 & Migración tabla transactions & Técnico 1 & Pendiente & & \\
\hline
HU-06 & Modelo Transaction con boot event & Técnico 1 & Pendiente & & \\
\hline
HU-06 & Componente BankMovementComponentForm & Técnico 1 & Pendiente & & \\
\hline
HU-06 & Pruebas unitarias HU-06 & Técnico 1 & Pendiente & & \\
\hline
HU-07 & Lógica año fiscal + lockForUpdate() & Técnico 1 & Pendiente & & \\
\hline
HU-07 & Pruebas unitarias HU-07 & Técnico 1 & Pendiente & & \\
\hline
HU-08 & Componente BankBookComponent & Técnico 2 & Pendiente & & \\
\hline
HU-08 & Consulta SQL con ventana SUM OVER & Técnico 2 & Pendiente & & \\
\hline
HU-08 & Pruebas de integración HU-08 & Técnico 2 & Pendiente & & \\
\hline
HU-10 & Migración tabla boxes & Técnico 1 & Pendiente & & \\
\hline
HU-10 & Modelo Box con boot event & Técnico 1 & Pendiente & & \\
\hline
HU-10 & Componente CashMovementComponentForm & Técnico 1 & Pendiente & & \\
\hline
HU-10 & Pruebas unitarias HU-10 & Técnico 1 & Pendiente & & \\
\hline
HU-11 & Lógica numeración mensual + lockForUpdate() & Técnico 1 & Pendiente & & \\
\hline
HU-11 & Pruebas unitarias HU-11 & Técnico 1 & Pendiente & & \\
\hline
HU-13 & Componente AccountStatementComponent & Técnico 2 & Pendiente & & \\
\hline
HU-13 & AccountStatementService con paginación & Técnico 2 & Pendiente & & \\
\hline
HU-13 & Pruebas de integración HU-13 & Técnico 2 & Pendiente & & \\
\hline
HU-15 & Componente ReportComponent & Técnico 2 & Pendiente & & \\
\hline
HU-15 & Exportación PDF con mPDF & Técnico 2 & Pendiente & & \\
\hline
HU-15 & Exportación Excel con Maatwebsite & Técnico 2 & Pendiente & & \\
\hline
HU-15 & Pruebas de exportación & Técnico 2 & Pendiente & & \\
\hline
\end{longtable}

\section{Preparación del Release}

El Sprint Final culmina con la Release Sprint Review y la Retrospectiva. A continuación se describen las actividades previas al release.

\subsection{Integración continua}

Todos los commits pasan por un pipeline de CI/CD que ejecuta las siguientes etapas:

\begin{enumerate}
    \item \textbf{Linting}: Ejecuta Pint (PHP CS Fixer) para verificar el estilo de código.
    \item \textbf{Pruebas}: Ejecuta Pest PHP con cobertura mínima del 80\%.
    \item \textbf{Compilación}: Compila assets con Vite (npm run build).
    \item \textbf{Despliegue}: Despliega automáticamente en el entorno de pre-producción tras pasar las pruebas.
\end{enumerate}

\subsection{Entorno de pre-producción}

Se despliega la versión completa del sistema en un entorno idéntico a producción (misma versión de PHP 8.3, MySQL 8.0, Nginx) para realizar pruebas finales de aceptación. Este entorno utiliza una copia anonimizada de la base de datos real.

\subsection{Pruebas de aceptación del Product Owner (PO)}

El PO (Gerente General) valida que todas las historias del Sprint Final cumplen los criterios de aceptación definidos en el Product Backlog:

\begin{itemize}
    \item El Auxiliar Contable puede registrar movimientos bancarios y de caja correctamente.
    \item El sistema asigna la numeración correlativa automática sin duplicados.
    \item El Contador puede consultar el libro de bancos y estados de cuenta con saldos actualizados.
    \item El Contador puede generar reportes con filtros y exportarlos a PDF y Excel.
    \item La interfaz es responsive y funciona correctamente en los navegadores objetivo.
\end{itemize}

\subsection{Capacitación a usuarios}

Se realizan talleres prácticos con los tres roles identificados:

\begin{itemize}
    \item \textbf{Auxiliar Contable (Carlos)}: Registro de movimientos de caja y bancos, generación de recibos.
    \item \textbf{Contador General (Carmen)}: Gestión de retenciones, consulta de estados de cuenta, generación de reportes.
    \item \textbf{Gerente (Pedro)}: Acceso al dashboard, consulta de reportes, supervisión general.
\end{itemize}

\subsection{Cierre del Sprint}

Se genera el \textbf{Incremento} del producto y se preparan las \textbf{Release Notes} con las siguientes funcionalidades completadas:

\begin{longtable}{|L{3cm}|L{5cm}|L{4cm}|}
\hline
\rowcolor[gray]{0.9}
\textbf{Módulo} & \textbf{Funcionalidades incluidas} & \textbf{Historias} \\
\hline
Administración & Usuarios, roles, permisos, proveedores, cuentas bancarias, impuestos & HU-17, HU-18, HU-19, HU-21, HU-22 \\
\hline
Retenciones & Registro, cálculo automático (RC-IVA, IUE, IT), numeración mensual, recibo PDF & HU-01, HU-02, HU-03, HU-04 \\
\hline
Libro de bancos & Movimientos bancarios, numeración fiscal, libro con saldo corriente & HU-06, HU-07, HU-08 \\
\hline
Caja & Movimientos de caja, numeración mensual & HU-10, HU-11 \\
\hline
Estados de cuenta & Consulta por proveedor/cliente, movimientos paginados, saldo corriente & HU-13 \\
\hline
Reportes & Reportes con filtros, exportación PDF y Excel & HU-15 \\
\hline
\end{longtable}

\newpage

\section{Avance de Interfaces Finales}

A continuación se muestra el estado de las interfaces principales del sistema al final del Sprint Final.

\begin{longtable}{|L{3.5cm}|L{2.5cm}|C{1.5cm}|C{1.5cm}|C{1.5cm}|C{1.5cm}|L{3cm}|}
\hline
\rowcolor[gray]{0.9}
\textbf{Módulo / Pantalla} & \textbf{Roles que la usan} & \textbf{Diseño UI} & \textbf{Frontend} & \textbf{Backend} & \textbf{Pruebas funcionales} & \textbf{Observaciones} \\
\hline
Login / Autenticación & Todos & $\checkmark$ & $\checkmark$ & $\checkmark$ & $\checkmark$ & Con animaciones y diseño responsivo \\
\hline
Dashboard (vista general) & Todos & $\checkmark$ & $\checkmark$ & $\checkmark$ & $\checkmark$ & KPIs básicos, retenciones y movimientos recientes \\
\hline
Lista de retenciones & Contador, Auxiliar & $\checkmark$ & $\checkmark$ & $\checkmark$ & $\checkmark$ & Filtro por mes y tipo (S/G) \\
\hline
Registro de retención & Contador & $\checkmark$ & $\checkmark$ & $\checkmark$ & $\checkmark$ & Cálculo automático de descuentos \\
\hline
Recibo PDF de retención & Contador & $\checkmark$ & $\checkmark$ & $\checkmark$ & $\checkmark$ & Formato A4, mPDF, monto en literal \\
\hline
Lista de movimientos bancarios & Auxiliar, Contador & $\checkmark$ & $\checkmark$ & $\checkmark$ & $\checkmark$ & Filtro por cuenta y mes \\
\hline
Registro movimiento bancario & Auxiliar & $\checkmark$ & $\checkmark$ & $\checkmark$ & $\checkmark$ & Con numeración fiscal automática \\
\hline
Libro de bancos (vista mensual) & Contador, Gerente & $\checkmark$ & $\checkmark$ & $\checkmark$ & $\checkmark$ & Saldo corriente con SUM OVER \\
\hline
Lista de movimientos de caja & Auxiliar, Contador & $\checkmark$ & $\checkmark$ & $\checkmark$ & $\checkmark$ & Filtro por mes \\
\hline
Registro movimiento de caja & Auxiliar & $\checkmark$ & $\checkmark$ & $\checkmark$ & $\checkmark$ & Con numeración mensual automática \\
\hline
Estados de cuenta & Contador, Gerente & $\checkmark$ & $\checkmark$ & $\checkmark$ & $\checkmark$ & Pestañas proveedor/cliente, saldo corriente \\
\hline
Reportes con filtros & Contador, Gerente & $\checkmark$ & $\checkmark$ & $\checkmark$ & $\checkmark$ & Exportación PDF y Excel \\
\hline
Gestión de usuarios & Administrador & $\checkmark$ & $\checkmark$ & $\checkmark$ & $\checkmark$ & CRUD completo con asignación de roles \\
\hline
Gestión de roles y permisos & Administrador & $\checkmark$ & $\checkmark$ & $\checkmark$ & $\checkmark$ & 35 permisos, sincronización con Spatie \\
\hline
Gestión de proveedores & Contador & $\checkmark$ & $\checkmark$ & $\checkmark$ & $\checkmark$ & Búsqueda por CI o nombre \\
\hline
Gestión de cuentas bancarias & Contador & $\checkmark$ & $\checkmark$ & $\checkmark$ & $\checkmark$ & Soporte BOB/USD/EUR \\
\hline
Gestión de impuestos/tasas & Contador & $\checkmark$ & $\checkmark$ & $\checkmark$ & $\checkmark$ & Parametrizable por tipo (S/G/A) \\
\hline
\end{longtable}

\newpage

\subsection{Diagrama de la interfaz de inicio de sesión}

La pantalla de inicio de sesión del sistema SIC-MINCH presenta un diseño moderno con fondo degradado y patrón de cuadrícula. Incluye el logo de la empresa MINCH SRL. y un formulario centrado con campos de correo electrónico y contraseña, más un botón de ingreso. Dispone de animaciones de entrada escalonadas para los campos del formulario y un efecto de brillo (glow) alrededor de la tarjeta de inicio de sesión.

\subsection{Diagrama del dashboard principal}

El dashboard principal se compone de cuatro secciones organizadas en un diseño de cuadrícula responsivo:

\begin{itemize}
    \item \textbf{Tarjetas KPI} en la parte superior: total de proveedores registrados, movimientos del mes, retenciones del mes, saldo de caja y saldo bancario consolidado.
    \item \textbf{Gráfico de tendencias} en el centro izquierdo: evolución de ingresos y egresos en los últimos 12 meses.
    \item \textbf{Tabla de movimientos recientes} en el centro derecho: últimos 10 movimientos registrados con fecha, descripción, tipo y monto.
    \item \textbf{Tabla de retenciones recientes} en la parte inferior: últimas 5 retenciones con código, proveedor y monto total.
\end{itemize}

\subsection{Diagrama del formulario de retención}

El formulario de registro de retención presenta los siguientes campos organizados en dos columnas:

\begin{itemize}
    \item Columna izquierda: tipo de retención (Servicios/Bienes) como selector, proveedor con buscador autocompletable, fecha con calendario.
    \item Columna derecha: monto base (input numérico con validación entre 10.00 y 999,999,999.99), descripción (área de texto).
    \item Sección inferior: tabla de descuentos calculados automáticamente al guardar, mostrando impuesto, tasa y monto.
    \item Botón de guardar con confirmación y mensaje de éxito con el código generado.
\end{itemize}

\subsection{Diagrama del libro de bancos}

La vista del libro de bancos presenta:

\begin{itemize}
    \item Selector de cuenta bancaria en la parte superior (dropdown con todas las cuentas registradas).
    \item Selector de mes/año para filtrar el período.
    \item Tabla con columnas: número de comprobante, fecha, descripción, débito, crédito, saldo corriente.
    \item El saldo corriente se calcula mediante una función de ventana SQL (SUM OVER) ordenada por fecha y número.
    \item Totales al final: suma de débitos, suma de créditos, saldo final.
    \item Botón para exportar la vista actual a PDF.
\end{itemize}

\subsection{Diagrama del estado de cuenta}

La vista de estados de cuenta presenta:

\begin{itemize}
    \item Pestañas en la parte superior para alternar entre \textbf{Proveedores} y \textbf{Clientes}.
    \item Lista de titulares con saldo actual (débito $-$ crédito) y un indicador visual de saldo positivo (verde) o negativo (rojo).
    \item Al seleccionar un titular, se muestra una tabla con todos sus movimientos: fecha, descripción, tipo, monto y saldo corriente.
    \item Paginación de movimientos (20 por página).
    \item Totales al final: suma débitos, suma créditos, saldo final.
    \item Botón para generar PDF del estado de cuenta completo.
\end{itemize}

\end{document}
